%================================
% Vocabulary & Language Focus — XB1 (G–L)
%================================
\documentclass[../main.tex]{subfiles}

\begin{document}

%================================
\section{G}
%================================

\vw{generate}{/ˈdʒenəreɪt/}{vt.}

\sense{vt.}{产生，引起}
\col{generate electricity / heat / power（物理的产生）}
\ex{The dam generates electricity for the whole valley.}
\col{generate interest / income / jobs / controversy（抽象的产生）}
\ex{The exhibition generated enormous interest among students.}
\ex{The new policy has generated a heated debate online.}

\begin{gram}{gen-“生”词族总动员}
拉丁语 \textit{generare}“生育”（\textit{genus}“族类”，原始印欧语 *\textit{genə-}“\textbf{生育}”）是英语的“生命之源”词根：\textbf{generate}（使生出$\to$产生）、\textbf{generation}（一代人；发电）、\textbf{generator}（发电机）、\textbf{regenerate}（re- 再生$\to$使再生；adj.\ 再生组织的）、\textbf{degenerate}（de- 离开族类$\to$堕落；退化）、\textbf{genesis}（创世；起源，\textit{Genesis}《创世记》）、\textbf{genetic}（基因的）、\textbf{genius}（天生之才，见下条）、\textbf{genuine}（天生的$\to$真正的）、\textbf{congenial}（同族的$\to$意气相投的）、\textbf{pregnant}（pre- 生育之前$\to$怀孕的“有孕育之兆的”）。凡见 \textit{gen-}，先想“生”。
\end{gram}

\vw{genius}{/ˈdʒiːniəs/}{n.}

\sense{n.}{天赋，天资}
\col{a genius for sth（对……有天赋，固定搭配 for）}
\ex{She has a genius for making complex ideas easy to understand.}
\col{a stroke of genius（神来之笔）}
\ex{Casting the comedy actor in the serious role was a stroke of genius.}

\sense{n.}{天才}
\col{a creative / mathematical / evil genius}
\ex{Einstein's genius changed physics forever.}
\ex{Mozart was a musical genius who composed from the age of five.}

\begin{origin}
拉丁语 \textit{genius}“\textbf{与生俱来的守护神}”，出自 \textit{gignere}“生育”（\textit{gen-} 词根，见上条）。罗马人相信每个男丁出生时都伴随一位\textbf{守护神灵}（genius），守护其性格与命运；它的女性对应者是 \textit{juno}。因守护神决定人的天性，genius 遂有“天生禀赋”义；文艺复兴时期又人格化为“天才人物”。“一出生就有的神”$\to$“天资”$\to$“天才”，三步走完两千年。同族的 \textbf{ingenious}（\textit{in-} 内在的天赋$\to$心灵手巧的；注意与 ingenuous“纯真的”一词之别）与阿拉伯故事借词 \textbf{jinn / genie}（灯神，拉丁 genius 的阿拉伯语化回借）都是它的亲戚。
\end{origin}

\vp{get down to doing sth}

\sense{v. 短语}{开始做某事，开始认真注意（或对待）某事}
\ex{Now that the celebrations are over, let's get down to work.}
\ex{It's time we got down to discussing the details of the plan.}

\begin{gram}{get 短语中的“to”：介词还是不定式？}
\textbf{get down to}、\textbf{get round to}（腾出手来做）、\textbf{get used to}、\textbf{look forward to}、\textbf{be devoted to}、\textbf{object to}、\textbf{be opposed to}——这些短语中的 \textbf{to 是介词}，后接\textbf{名词或动名词}：\ding{55} \textit{get down to discuss}；\checkmark\ \textit{get down to discussing}。判断技巧：to 后能直接跟名词（\textit{look forward to the holiday}）即是介词。语法填空中此类“伪装不定式”年年设卡。
\end{gram}

\vw{gifted}{/ˈɡɪftɪd/}{adj.}

\sense{adj.}{有才华的，有天赋的}
\col{a gifted musician / player / child}
\ex{The school runs special programs for gifted children.}
\ex{Even gifted athletes need years of disciplined training.}

\begin{gram}{“天赋”名词化：gift 的隐转喻}
\textbf{gift} 本义“赠品”；中世纪起“上天所赠之才”义渐显（\textit{the gift of the gab} 天生的好口才；\textit{have a gift for languages} 有语言天赋）。同构的隐转喻在多种语言中平行：德语 \textbf{Begabung}（天赋$<$geben 给）、法语 \textbf{don}（才能$<$ donner 给）——“天赋是神赐的礼物”是跨文化的共同隐喻。注意 \textbf{gifted} 与 \textbf{talented} 的语体差：gifted 多用于艺术与学术，talented 泛指才能，语体更平。
\end{gram}

\vw{glorious}{/ˈɡlɔːriəs/}{adj.}

\sense{adj.}{荣耀的，光荣的}
\col{a glorious victory / history / tradition}
\ex{The museum recalls the city's glorious past as a trading port.}

\sense{adj.}{壮丽的，辉煌的}
\col{glorious scenery / sunshine / colors}
\ex{We had a glorious week of sunshine in the mountains.}
\col{glory（名词：光荣；壮美）}
\ex{The team brought glory to their country.}

\begin{origin}
经由古法语 \textit{glorieus} $\leftarrow$ 拉丁语 \textit{gloriosus}“\textbf{充满荣耀的}”，出自 \textit{gloria}“荣耀、声誉”（词源成谜：一说借自日耳曼语 *\textit{klēraz}（与英语 \textbf{loud} 同根，“被高声传颂”$\to$名声），一说与 \textit{clue/cla} 无关的土著拉丁词）。\textit{Gloria in excelsis Deo}（荣耀归于至高的神）使它染上宗教光辉；18 世纪后英语中发展出反讽用法 \textit{a glorious mess}“一团堂皇的糟”——\textbf{gloriously} 与失败连用（\textit{gloriously wrong} 错得离谱又坦然）已成为英式幽默的固定配方。
\end{origin}

\vp{go about sth}

\sense{v. 短语}{忙于做某事，继续做某事；着手处理}
\ex{Despite the rumors, she went about her work as usual.}
\ex{This is how professionals go about solving such problems.}
\col{go about your business（做自己的事）}
\ex{The guards told us to go about our business and move along.}

\vw{grab}{/ɡræb/}{vt. \& vi.}

\sense{vt. \& vi.}{抓住；赶紧，抓紧}
\col{grab sth / grab at sth（抓住／试图抓住）}
\ex{He grabbed my arm and pulled me back from the road.}
\ex{She grabbed at the railing but missed.}
\col{grab a bite / a coffee / a taxi（口语：赶快弄到）}
\ex{Let's grab a quick bite before the meeting.}
\ex{I grabbed a taxi outside the hotel.}

\sense{vt.}{吸引（注意力）}
\col{grab sb's attention（$=$ catch sb's attention）}
\ex{The headline was designed to grab readers' attention.}

\begin{gram}{“抓取”动词的力度谱系}
\textbf{take}（中性拿）$<$ \textbf{grab}（猛地一把，常含匆忙或粗鲁）$<$ \textbf{seize}（强力夺占，书面；\textit{seize an opportunity} 抓住时机是褒义习语）$<$ \textbf{snatch}（出其不意地抢夺）$<$ \textbf{grasp}（紧握；理解，见下条）$<$ \textbf{grip/clutch}（死死攥住；\textit{clutch} 常含惊慌）。写作中“抓机会”用 \textbf{seize} 或 \textbf{grasp}，勿用 grab（显粗鲁）；“吸引注意”则 grab 与 attract / draw 均可。
\end{gram}

\vw{grand}{/ɡrænd/}{adj.}

\sense{adj.}{宏大的，宏伟的；壮丽的，堂皇的}
\col{a grand building / palace / ceremony}
\ex{The palace's grand staircase is lined with marble statues.}
\col{grand plans / ambitions / designs（宏图）}
\ex{The young CEO was full of grand plans for expansion.}
\col{a grand total of（总计，戏谑）}
\ex{All my savings amounted to a grand total of 87 yuan.}

\begin{origin}
经由古法语 \textit{grant}“伟大的” $\leftarrow$ 拉丁语 \textit{grandis}“\textbf{巨大的、年长的、充盈的}”（原始印欧语一说 *\textit{ghreu-}“燃烧旺盛”？词源不确，可能与“生长成熟”相关）。法语中 \textit{grand} 一词拼法未变，但在英语里被 15 世纪的抄写员凭空添上一个 \textbf{-d}（原本只拼 \textit{graund}），\textbf{d 不发音}——与 debt 的哑 b 同属“字母镀金”。它的后代枝繁叶茂：\textbf{grandmother}（祖母“更长辈的母”）、\textbf{grand piano}（大钢琴）、\textbf{grandeur}（壮丽）、\textbf{aggrandize}（ag- 使大$\to$吹捧扩张）；而 \textbf{grand} 美俚“一千美元”则出自 \textit{a grand}＝\textit{a thousand dollars}的窃贼黑话。
\end{origin}

\vw{grasp}{/ɡrɑːsp/}{vt.\ \&\ n.}

\sense{vt.}{抓紧；领会，理解}
\col{grasp sth（本义握紧）}
\ex{She grasped the rope with both hands and pulled hard.}
\col{grasp an idea / a concept / the significance（领会）}
\ex{The students quickly grasped the basic idea of the theory.}
\col{grasp an opportunity（抓住时机）}
\ex{He grasped the chance to study abroad when it came.}

\sense{n.}{理解，领会；抓紧，握紧，控制；能力所及}
\col{have a good grasp of sth（对……透彻掌握）}
\ex{A good grasp of grammar is essential for translators.}
\col{beyond one's grasp / within one's grasp（力所不及／唾手可得）}
\ex{The championship is now within their grasp.}
\ex{Such abstract ideas are beyond the grasp of young children.}

\begin{gram}{“理解”的隐喻地图：握、抓、逮、吞}
英语表“理解”的动词几乎全是身体动作的隐喻：\textbf{grasp}（握住）、\textbf{catch}（逮住：\textit{Do you catch my meaning?}）、\textbf{get}（得到：\textit{I get it.}）、\textbf{follow}（跟随：\textit{Can you follow me?}）、\textbf{digest}（消化，见 D 部分）、\textbf{take in}（吞入：\textit{That's a lot to take in.}）。汉语同样说“\textbf{把握}”“\textbf{吃透}”——身体经验（抓住才可靠、吃下才属于自己）塑造了各语言共通的认知隐喻，这是认知语言学（Lakoff \& Johnson）的奠基观察。
\end{gram}

\vw{grave}{/ɡreɪv/}{n.\ \&\ adj.}

\sense{n.}{坟墓；死亡，去世}
\col{from the cradle to the grave（从摇篮到坟墓，一生）}
\ex{The welfare system follows a citizen from the cradle to the grave.}
\col{dig one's own grave（自掘坟墓）}
\ex{By lying to the committee, he dug his own political grave.}
\col{turn in one's grave（死不瞑目的愤慨）}
\ex{The founder would turn in his grave if he saw the factory now.}

\sense{adj.}{严重的；严肃的}
\col{a grave mistake / danger / situation（严重）}
\ex{The doctor's silence suggested something grave.}
\col{a grave expression / voice（严肃）}
\ex{``Sit down,'' he said in a grave voice.}
\col{gravity（名词：严重性；重力——一词两义同源）}
\ex{They failed to grasp the gravity of the situation.}

\begin{origin}
两个 grave 是\textbf{词源各异}的同形词。名词 grave（坟）是古英语 \textit{græf}“\textbf{挖掘之物}”，与动词 \textbf{dig} 的古词 \textit{grafan} 同根（原始印欧语 *\textit{ghrebh-}“挖、刮”，同族 \textbf{engrave}“雕刻”“挖进去”，德语 \textbf{Grab} 坟）。形容词 grave（严重的）经由法语 $\leftarrow$ 拉丁语 \textit{gravis}“\textbf{沉重的}”（原始印欧语 *\textit{gwerə-}“重”，同根 \textbf{grief} 悲伤、\textbf{aggravate} 加重、\textbf{gravid} 怀孕的“负重的”、\textbf{heavy}?（否——heavy 是日耳曼本地词，仅义近））。故 \textbf{gravity} 兼指“重力”与“严重性”——牛顿与伦理学家共用一个“沉重”。
\end{origin}

\vw{greedily}{/ˈɡriːdɪli/}{adv.}

\sense{adv.}{贪婪地，贪心地}
\ex{The children tore open the parcels and greedily devoured the sweets.}
\col{greedy / greed（词族）}
\ex{He was greedy for power and fame.（greedy for sth）}
\ex{Her greed for land ruined the family.}

\begin{gram}{形容词 $+$ -ly 副词化的拼式陷阱}
\textbf{greedy $\to$ greedily}（y 变 i 加 -ly），同类：\textit{easy $\to$ easily, happy $\to$ happily, heavy $\to$ heavily}。注意 \textbf{greed} 是名词（无 y）；\textbf{greediness} 是名词的名词化，语体繁复，慎用。七宗罪（the seven deadly sins）英语名录值得顺带记诵：\textbf{greed}（贪婪）、\textbf{pride}（傲慢）、\textbf{envy}（嫉妒）、\textbf{wrath}（暴怒）、\textbf{gluttony}（暴食）、\textbf{lust}（色欲）、\textbf{sloth}（懒惰）——宗教语篇与文学评论的常客。
\end{gram}

\vw{greet}{/ɡriːt/}{vt.}

\sense{vt.}{问候，欢迎，招呼}
\col{greet sb with a smile / a nod}
\ex{She greeted us with a warm smile and a pot of tea.}

\begin{gram}{be greeted with / by：事物“迎来”何种反应}
greet 的被动式是新闻英语的高频句式，主语常为\textbf{事物}：“\textit{The announcement was greeted \textbf{with} silence.} 公告迎来的是沉默。”\textit{be greeted by applause / cheers / suspicion / relief}。介词 with 接反应，by 接施动者。写作中用它替代 \textit{people reacted to\dots} 更显地道。
\end{gram}

\vw{guarantee}{/ˌɡærənˈtiː/}{vt.\ \&\ n.}

\sense{vt.}{保证，担保；确保；为（产品）提供保修}
\col{guarantee that\dots / guarantee to do sth}
\ex{I guarantee that the repairs will be finished by Friday.}
\ex{No one can guarantee to arrive on time in such weather.}
\col{guarantee sth against sth（为……提供担保）}
\ex{The camera is guaranteed against water damage for two years.}

\sense{n.}{保证，担保；保修单}
\col{under guarantee（在保修期内）}
\ex{The laptop is still under guarantee, so the repair is free.}
\col{a guarantee of / that\dots}
\ex{Wealth is no guarantee of happiness.}

\begin{origin}
它是英语中最老的“法律德语借词”之一：16 世纪直接借自德语 \textbf{Garantie}／\textit{garantieren}，溯其源则经西班牙语 \textit{garantía}（担保人）$\leftarrow$ 晚期拉丁语/西哥特语境的 *\textit{warandī}“\textbf{郑重许诺}”，与英语本土词 \textbf{warrant}（保证；逮捕令；担保）、\textbf{warranty}（保修）是\textbf{同一个日耳曼语词}经罗曼语“出口转内销”的双胞胎——warrant 是法语带进来的近亲，guarantee 是德语带回来的另一支。拼写上 gu- 读 /ɡ/（guard, guess, guitar 同例），是诺曼抄写员防止 /g/ 被误读 /j/ 加的 u。一词根、两条商路，写着中世纪地中海的贸易史。
\end{origin}

%================================
\section{H}
%================================

\vw{heaven}{/ˈhevn/}{n.}

\sense{n.}{极乐之地}
\col{heaven on earth（人间天堂）}
\ex{The quiet island felt like heaven on earth to the tired travelers.}
\col{be in heaven（乐上天了，口语）}
\ex{A hot bath and a good book --- I was in heaven.}
\col{move heaven and earth（竭尽全力）}
\ex{They moved heaven and earth to get tickets for the final.}
\col{the heavens（书面/古雅：天空）}
\ex{Galileo turned his telescope to the heavens.}

\begin{origin}
古英语 \textit{heofon}“\textbf{天空、天界}”（原始日耳曼语 *\textit{hibin-}，词源诸说纷纭：一说与“\textbf{举起、覆盖}”相关，即“覆于头顶之幕”）。宗教语义分层值得注意：单数 \textbf{heaven} 在基督教语境是上帝居所；复数 \textbf{the heavens} 保留古义“苍穹”（天文语篇仍用：\textit{the starry heavens}）；俚语感叹 \textit{Heavens!} / \textit{Good heavens!} 则是维多利亚时代为避讳 God 而生的委婉语。汉语“天”兼指天空、天道与上帝，与 heaven 的语义分层惊人平行——单音节汉字与多义日耳曼词在此遥相印证。
\end{origin}

\vw{hence}{/hens/}{adv.}

\sense{adv.}{因此，由此（正式书面语）}
\ex{The town lies at the meeting of two rivers; hence its importance to trade.}
\ex{The data was incomplete; hence the delay in publishing the report.}

\begin{gram}{hence / therefore / thus / consequently 的语体阶梯}
\textbf{therefore} 通用正式，前后需较完整句法；\textbf{hence} 最文言，\textbf{后面可直接跟名词短语}（\textit{hence the delay}），这是它独有的“悬垂用法”；\textbf{thus} 可接分词（\textit{thus ending the dispute}）；\textbf{consequently} 强因果链条，多用于说明文。学术写作中 hence 常与公式、定义连用（\textit{hence the theorem}）。另注意古用法 \textit{two weeks hence}（从今两周后，“由此起算”）仍见于法律文书。
\end{gram}

\begin{origin}
中古英语 \textit{hennes}“\textbf{从这里}”：\textit{henne}（此地）+ 表示方向的属格后缀 \textit{-s}——与 \textbf{hence}“由此”、\textbf{thence}“由彼”、\textbf{whence}“由何”构成一套完整的中古英语\textbf{方位指示系统}：\textit{hither/hence/here}（近指）、\textit{thither/thence/there}（远指）、\textit{whither/whence/where}（疑问）。现代英语只保留了 here/there/where 一列，其余皆成化石词，唯有 hence 因逻辑连接功能在学术语体中复活。“从这里”$\to$“从这一点（推出）”，空间指示词语法化为逻辑连接词，是语法化的经典案例。
\end{origin}

\vw{historic}{/hɪˈstɒrɪk/}{adj.}

\sense{adj.}{历史上著名的；有史时期的}
\col{a historic moment / decision / victory（载入史册的）}
\ex{The two leaders shook hands in a historic first meeting.}
\col{in historic times（有文字记载以来，与 prehistoric 相对）}
\ex{The volcano last erupted in historic times, in 1660.}

\begin{lex}{historic / historical}
\textbf{historic}：“\textbf{有历史意义的}、载入史册的”（\textit{a historic summit}）；\textbf{historical}：“\textbf{与历史相关的}、历史上存在过的”（\textit{historical documents / novels / figures}）。一个事件可以 historical 而毫不 historic（平凡的历史旧事）。口诀：\textbf{historic 留名，historical 沾边}。另注意 \textit{Historic Present}（历史现在时）——叙述往事用现在时以增临场感（\textit{In 1492 Columbus sails west\dots}），是文学修辞一枝。
\end{lex}

\begin{origin}
经由拉丁语 \textit{historicus} $\leftarrow$ 希腊语 \textit{historikos}，出自 \textit{historia}“\textbf{探究、见闻、记事}”（动词 \textit{historein}“探问、考证”，源自 \textit{histōr}“\textbf{智者、见证人}”，原始印欧语 *\textit{weid-}“看见、知道”——与 \textbf{wit}（机智）、\textbf{wise}、\textbf{vision}（\textit{vidēre} 看见）、\textbf{idea}（\textit{idein} 看到$<$\textit{wid-}）、\textbf{evident}（看得明的）同根）。“历史”的本义不是过去，而是\textbf{亲眼见证后的探究}——希罗多德以此命名其著作，遂成学科之名。知与看同根，故英语说 \textit{I see} 表“我懂”。
\end{origin}

\vp{hit the spot}

\sense{v. 短语}{使人满意，正合需要（多指饮食）}
\ex{A cold lemonade on a hot afternoon really hits the spot.}
\ex{That bowl of noodles hit the spot after the long hike.}

\begin{gram}{“spot”短语小家族}
\textbf{hit the spot}（正中所需，原为广告语“一击即中要穴”）；\textbf{on the spot}（当场；立即：\textit{He answered the question on the spot.}；\textit{put sb on the spot} 使某人当场为难）；\textbf{spot check}（抽查）；\textbf{a soft spot for sb/sth}（偏爱：\textit{I have a soft spot for old movies.}）；\textbf{spot sb doing sth}（当场发现：\textit{The guard spotted him climbing the wall.}）——spot 由“点”到“现场”到“发现”，一条词义链。
\end{gram}

%================================
\section{I}
%================================

\vw{ideal}{/aɪˈdiːəl/}{n.\ \&\ adj.}

\sense{n.}{理想；典范}
\col{live up to one's ideals（践行理想）}
\ex{He found it hard to live up to his own high ideals.}
\col{the ideal of liberty / beauty（……之理想）}
\ex{The statue embodies the Greek ideal of physical perfection.}

\sense{adj.}{完美的，理想的}
\col{be ideal for sth}
\ex{The weather was ideal for an outdoor wedding.}
\ex{This location is ideal for a small café.}

\begin{origin}
源自晚期拉丁语 \textit{idealis}“\textbf{属于理念（idea）的}”，而 \textbf{idea} 经由拉丁语 $\leftarrow$ 希腊语 \textit{idea}“\textbf{形、模样、理念}”，出自 \textit{idein}“\textbf{看见}”（原始印欧语 *\textit{weid-}“见、知”，与 historic、vision、wit 同根，见上条）。柏拉图据此造出\textbf{理念论}（Theory of Forms）：眼见之物皆是“理念”（\textit{idea}）的 imperfect 摹本——故 \textbf{ideal} 生来就是“\textbf{彼岸的完美范型}”：世上的理想主义者（idealist），按词源，都是柏拉图的门徒。同根的 \textbf{idol}（偶像）$<$希腊 \textit{eidōlon}“影像、幻象”是它的孪生兄弟——崇拜的偶像与追求的理想，在词源上一个是“假象”，一个是“真形”。
\end{origin}

\vw{imply}{/ɪmˈplaɪ/}{vt.}

\sense{vt.}{含有……的意思，暗示，暗指}
\col{imply that\dots（暗指）}
\ex{His silence implied agreement, though he never said yes.}
\col{be implied in sth（隐含于）}
\ex{Risk is implied in any investment.}
\col{implication（名词：含意；牵连；影响）}
\ex{What are the implications of this decision for small businesses?}
\ex{He was arrested for his alleged implication in the plot.}

\begin{gram}{imply / infer：说话人与听话人的分工}
\textbf{imply} 是\textbf{说话人}暗含（\textit{The writer implies that\dots}）；\textbf{infer} 是\textbf{听话人}推断（\textit{The reader infers that\dots}）。经典记忆句：\textit{The speaker implies, the listener infers.} 考试与编辑实务中两者常被混用（\ding{55} \textit{Are you implying I'm wrong?} 口语可，但 \textit{I imply from your words\dots} 是错的，应为 infer）。完形与阅读设问 “What can be inferred\dots?” 正是考 infer 的这一“读者视角”。
\end{gram}

\begin{origin}
经由拉丁语 \textit{implicare}“\textbf{折入、编入}”：\textit{im-}（入）+ \textit{plicare}“折叠”（\textit{plec-}“编”词根的近亲，见 complex、display、employ 词条）。“暗示”即把意思\textbf{折叠进}话语里，听者须自行展开——fold in 与 unfold（\textbf{explicate} 说明，ex- 展开之）构成一对镜像动词：imply 折入、explicate 展开，故 \textbf{implicit}（暗含的）与 \textbf{explicit}（明白的，ex-plicare 展开）是一对标准反义词。同根的还有 \textbf{implicate}（把人\textbf{卷折进}案件$\to$牵连）。
\end{origin}

\vw{impress}{/ɪmˈpres/}{vt. \& vi.}

\sense{vt. \& vi.}{使钦佩，给……留下深刻印象}
\col{impress sb (with sth) / be impressed by/with}
\ex{The young pianist impressed the judges with her composure.}
\ex{We were deeply impressed by the volunteers' devotion.}
\col{impress sth on/upon sb（使铭记）}
\ex{My teacher impressed upon me the value of honesty.}

\sense{vt.}{使意识到（重要性等）}
\ex{The report impressed upon the committee the urgency of reform.}
\col{be impressed into service（被强征）}
\ex{Fishing boats were impressed into naval service during the war.}

\begin{gram}{press 词族回访与 impress 的两个被动}
\textbf{impress} $=$ im-（入）+ press（压，词族见 depress 词条）。“压入”的三次引申：（1）印章压入蜡泥$\to$\textbf{铭刻}（\textit{impress one's seal}）；使铭记$\to$印象深刻（\textbf{impression} 印象；\textbf{impressive} 令人钦佩的）；（2）强征入伍（\textit{impress sailors}，历史术语“抓壮丁”）；（3）印刷术语“压印”（\textit{imprint} 版记）。第一人称与第三人称的反应不同词：事物 \textbf{is impressive}，人 \textbf{is impressed}——写作中勿混。
\end{gram}

\vp{in one's own right}

\sense{adv. 短语}{凭自身的资格（或努力）}
\ex{She is a famous singer in her own right, not just the daughter of a star.}
\ex{The novel is a classic in its own right, apart from its author's fame.}

\begin{gram}{right 的“权利”词族与习语}
\textbf{in one's own right} 中 right 取“法定资格”义：因\textbf{自身}的名分（而非依附他人）而成立。相关短语：\textbf{by right(s)}（按理说：\textit{By rights, the prize is hers.}）、\textbf{within one's rights}（有权如此）、\textbf{rightly or wrongly}（不管对错）。\textbf{right} 的两义（右；权利）在法语 \textit{droit}、德语 \textit{Recht} 中同样一字双关——右手为正为权，是各文化共有的身体政治学。
\end{gram}

\vw{industrialization}{/ɪnˌdʌstriəlaɪˈzeɪʃn/}{n.}

\sense{n.}{工业化}
\col{rapid / large-scale industrialization}
\ex{The region's rapid industrialization transformed it within a single generation.}
\col{词族：industry / industrial / industrialized}
\ex{The car industry is the backbone of the local economy.}
\ex{Japan and Germany are highly industrialized economies.}

\begin{origin}
\textbf{industry} 来自拉丁语 \textit{industria}“\textbf{勤勉}”（形容词 \textit{industrius}“勤奋的”），析为 \textit{indu-}（在内）+ \textit{struere}“\textbf{建造}”（\textit{stru-}“堆积建造”词根：\textbf{structure} 结构、\textbf{construct} 建造、\textbf{instruct}（在体内筑入$\to$教导）、\textbf{instrument} 工具、\textbf{destroy}（de- 拆建$\to$毁坏））——本义“\textbf{在内里构筑}”即“勤勉”。18 世纪，“勤勉品质”转指“勤勉的集体行为”$\to$\textbf{产业}（the cotton industry），再机器化$\to$\textbf{工业}；形容词 industrial、动词 industrialize、名词 industrialization 逐级派生。一个“勤奋”的人格词长成了一部\textbf{工业革命}的编年史。
\end{origin}

\vw{ingredient}{/ɪnˈɡriːdiənt/}{n.}

\sense{n.}{材料，成分；因素，要素}
\col{ingredients for sth（……的用料，注意 for）}
\ex{Buy the ingredients for the cake on your way home.}
\col{the ingredients of / for success（成功的要素，喻义高频）}
\ex{Hard work and luck are the main ingredients of his success.}
\ex{Tolerance is an essential ingredient in any friendship.}

\begin{origin}
拉丁语 \textit{ingrediens}（\textit{ingredi}“\textbf{进入}”的现在分词）：\textit{in-}（入）+ \textit{gradi}“\textbf{迈步}”（\textit{gradi} 是 \textit{gradus}“一步、台阶”的动词，原始印欧语 *\textit{ghredh-}“行走”）。“配料”的本义是“\textbf{走进去（锅里/面团）的东西}”。同根词的“步伐”家族：\textbf{grade}（级别“台阶”）、\textbf{gradual}（一步步的）、\textbf{graduate}（修完台阶$\to$毕业；刻度）、\textbf{aggressive}（ag- 向前迈$\to$进取的；挑衅的）、\textbf{progress}（向前走$\to$进步）、\textbf{congress}（走到一起$\to$代表大会）、\textbf{transgress}（越步而过$\to$逾越）。食谱与议会在同一个“迈步”词根上会师。
\end{origin}

\vw{inner}{/ˈɪnər/}{adj.}

\sense{adj.}{内心的，隐藏的；里面的，内部的}
\col{inner thoughts / feelings / peace（内心世界）}
\ex{Yoga helps her maintain a sense of inner peace.}
\col{the inner circle（核心圈子）}
\ex{Only the inner circle of advisers knew the plan.}
\col{an inner city（市中心贫民区）；inner ear（内耳）}
\ex{The program funds music classes in inner-city schools.}

\begin{gram}{inner / inward / internal / interior}
\textbf{inner} 兼指物理内部与内心（最常用于心理）；\textbf{inward} 方向性“向内的”（\textit{an inward glance} 内省；\textit{inward investment} 内部投资）；\textbf{internal} 正式，指机构/身体/国家之内（\textit{internal affairs} 内政；\textit{internal injuries} 内伤）；\textbf{interior} 空间/地理内部（\textit{interior design} 室内设计；\textit{the interior} 内陆腹地）。反义阵营：outer / outward / external / exterior。选词先辨“身体的、机构的、地理的还是心灵的”。
\end{gram}

\vw{innovation}{/ˌɪnəˈveɪʃn/}{n.}

\sense{n.}{创造，创新；新思想，新方法}
\col{technological / scientific innovation}
\ex{The country's future depends on continuous technological innovation.}
\col{an innovation in sth（……上的革新）}
\ex{The smartphone was a remarkable innovation in communication.}
\col{innovate / innovative（动词／形容词）}
\ex{The company encourages staff to innovate.}
\ex{an innovative approach to teaching}

\begin{gram}{nov-“新”词族}
拉丁语 \textit{novus}“新”（原始印欧语 *\textit{newo-}，与英语 \textbf{new} 同源——nov- 是拉丁形，new 是日耳曼形）：\textbf{innovate}（in- 引入新$\to$革新）、\textbf{renovate}（re- 再翻新$\to$翻修）、\textbf{novelty}（新奇；\textbf{novel} 小说“新奇之事”）、\textbf{novice}（新手上路$\to$ novice teacher）、\textbf{novelty shop} 新奇商品店。写作模板：\textit{Innovation distinguishes a leader from a follower.}（乔布斯语的平实版），\textit{innovative solutions to old problems} 是议论文万能短语。
\end{gram}

\vw{insight}{/ˈɪnsaɪt/}{n.}

\sense{n.}{洞悉，了解；洞察力，领悟}
\col{an insight into sth（对……的洞见，固定介词 into）}
\ex{The documentary offers a rare insight into village life in the 1930s.}
\ex{Her speech gave us an insight into how a champion thinks.}
\col{gain / provide / have an insight}
\ex{Reading her diary, historians gained fresh insights into the court.}
\col{insightful（富有洞察力的）}
\ex{It is an insightful analysis of the crisis.}

\begin{origin}
\textbf{in-}（入）+ \textbf{sight}（看见）——“\textbf{看进去}”的能力。sight 的词根 *\textit{sekwh-}“跟随目光”？实为 \textit{sekh\textsuperscript{w}-}；其拉丁表亲 \textit{specere/spicere}“看”更配得上“洞察”家谱：\textbf{inspect}（往里看$\to$检查，见下条）、\textbf{respect}（一再回看$\to$尊重）、\textbf{suspect}（由下看$\to$怀疑）、\textbf{perspective}（透过去看$\to$透视；观点）、\textbf{retrospect}（回头看$\to$回顾）、\textbf{prospect}（向前看$\to$前景）。英语说“洞见”靠\textbf{目光的路线}（in-/re-/pro-），汉语说“洞见”靠\textbf{穿墙的目光}（洞）——两种语言各自给“理解”画了一张视觉地图。
\end{origin}

\vw{inspection}{/ɪnˈspekʃn/}{n.}

\sense{n.}{检查，查看，审视；视察}
\col{an inspection of sth / carry out an inspection}
\ex{Officers carried out a routine inspection of the factory.}
\col{on closer inspection（细看之下，叙事高频）}
\ex{The vase looked genuine, but on closer inspection it was a modern copy.}
\col{inspect / inspector}
\ex{Officials inspect the bridge every two years.}
\ex{A ticket inspector walked through the carriage.}

\vw{interpret}{/ɪnˈtɜːprɪt/}{vt.\ \&\ vi.}

\sense{vt.}{解释，阐释；把……理解为，领会}
\col{interpret sth as sth（把……理解为）}
\ex{His refusal was interpreted as a lack of interest.}
\col{interpret a dream / the law / data（解读）}
\ex{Two courts interpreted the same law in opposite ways.}

\sense{vi.}{口译}
\col{interpret for sb（为某人担任口译）}
\ex{She interpreted for the visiting delegation throughout the week.}

\begin{gram}{interpret / translate：口译与笔译分工}
\textbf{translate} 笔译（书面）；\textbf{interpret} 口译（现场）——名词 \textbf{interpreter}（口译员；也可能是“解说者”）与 \textbf{translator}（笔译员）分工对应。注意：interpret 的“口译”义\textbf{不用于}“把 A 语言译成 B 语言”的及物句型（\ding{55} \textit{interpret the speech into English}），要么 \textit{interpret for sb}，要么用 \textit{translate\dots into}。同根名词 \textbf{interpretation} 兼“解释、演绎、口译”三义（\textit{a bold interpretation of the sonata} 对奏鸣曲的大胆演绎）。
\end{gram}

\begin{origin}
经由拉丁语 \textit{interpretari}“\textbf{解释、居间传话}”，出自 \textit{interpres}“\textbf{中间人、经纪、译员}”（\textit{inter-}“之间”+ 词根 \textit{-pres}，或与买卖中介相关）。词源核心是\textbf{居间}：译员站在两种语言\textbf{之间}，阐释者站在文本与读者\textbf{之间}。衍生词 \textbf{interpretation} 在 19 世纪哲学中大放异彩——施莱尔马赫与狄尔马的诠释学（\textbf{hermeneutics}，另出希腊 \textit{hermēneuein}，与赫尔墨斯 Hermes“神界信使”相关：众神的\textbf{interpreter}）正是研究“居间理解”的学问。
\end{origin}

\vw{invader}{/ɪnˈveɪdər/}{n.}

\sense{n.}{武装入侵的军队（或国家），侵略者}
\col{repel / drive out an invader}
\ex{The villagers resisted the invaders with farm tools.}
\col{invade / invasion（词族）}
\ex{The army invaded the border region at dawn.}
\ex{The invasion was launched without warning.}

\begin{origin}
\textbf{invade} $\leftarrow$ 拉丁语 \textit{invadere}“\textbf{走进去、攻入}”：\textit{in-}（入）+ \textit{vadere}“\textbf{行进}”（原始印欧语 *\textit{wadh-}“走”，同根 \textbf{pervade}（per- 遍走$\to$弥漫）、\textbf{evade}（e- 走脱$\to$规避）、\textbf{invade}（走入$\to$入侵）、\textbf{wade}（涉水而行，日耳曼表亲））。“侵略”的词源只是“\textbf{走进去}”——词汇的中性原料（走路）加军事语境炼成刀剑。有趣的是 \textit{vadere} 还给了经济学家 \textbf{asset}？否——asset 出于法语 \textit{assez}“足够”；但 \textbf{vade mecum}（随身手册“与我同行”，拉丁语短语）至今是学术出版的小众书名。
\end{origin}

%================================
\section{J}
%================================

\vw{jazz}{/dʒæz/}{n.}

\sense{n.}{爵士乐}
\ex{Louis Armstrong brought jazz from New Orleans to the world.}
\col{jazz sth up（使更有生气，口语）}
\ex{They jazzed up the old song with electronic beats.}
\col{jazz hands（爵士手，音乐剧手势）}

\begin{origin}
词源是英语词汇史最大的悬案之一：首见于 1912 年加州报纸（原拼 \textit{jasm}，最初竟是\textbf{棒球术语}“活力、爆发力”，指投手的球路），1915 年前后移用于芝加哥的新音乐。候选来源包括：方言 \textit{jasm}“精力”（$<$\textit{jism}）、非洲某语言、克里奥尔语。学术考证的最新共识倾向于“俚语活力义先行，乐名随后”。与 blues（见 B 部分）合观：20 世纪上半叶美国黑人音乐几乎为英语贡献了整套词汇——\textbf{rock and roll}（1950s）、\textbf{blues}、\textbf{jazz}、\textbf{gospel}（福音），它们的词源多不可考，恰似其音乐：即兴，无名，而生猛。
\end{origin}

%================================
\section{L}
%================================

\vw{landscape}{/ˈlændskeɪp/}{n.}

\sense{n.}{乡村风景画；风景，景色}
\col{a bleak / dramatic / urban landscape}
\ex{The train wound through a landscape of lakes and pine forests.}
\col{the political / economic landscape（喻：格局）}
\ex{The scandal changed the political landscape overnight.}
\col{landscape gardening（园林营造）}
\ex{He studied landscape gardening in Britain.}

\begin{gram}{-scape 系列与“化零为整”后缀}
\textbf{landscape} $<$ 荷兰语 \textit{landschap}“地区画”（\textit{land}+\textit{schap}“-ship 状态/技艺”），16 世纪作为\textbf{绘画术语}进入英语——先有“风景画”，后有“风景”本身！后缀 -scape $=$ \textbf{-ship}（friendship, hardship）的荷兰拼法。类比造词由此生发：\textbf{seascape}（海景画）、\textbf{cityscape}（城市景观）、\textbf{moonscape}（月面荒景）、\textbf{dreamscape}（梦境景象）、乃至数字时代的 \textbf{mediascape / cyberscape}。掌握“-scape $=$ ……的全景样貌”，科技文阅读中的生词自解。
\end{gram}

\vw{lane}{/leɪn/}{n.}

\sense{n.}{（乡间）小路；小巷；车道；跑道；航线}
\col{a country lane（乡间小路）}
\ex{We cycled along quiet country lanes all morning.}
\col{a bus / bicycle lane（车道）}
\ex{Cars are not allowed in the bus lane during rush hours.}
\col{the fast / slow lane（快／慢车道；喻人生节奏）}
\ex{He enjoyed life in the slow lane after retirement.}
\col{ocean lanes（航线）}
\ex{Ships follow established ocean lanes to avoid collisions.}

\vw{layer}{/ˈleɪər/}{n.}

\sense{n.}{层，表层；层次}
\col{a layer of sth（一层……）}
\ex{Spread a layer of cheese over the dough.}
\ex{A thick layer of snow covered the fields.}
\col{the layers of meaning（层层含义）}
\ex{The poem rewards rereading with its many layers of meaning.}
\col{layered（分层的；层次丰富的）}
\ex{The dessert is made of three flavored layers, simply layered and chilled.}

\begin{gram}{lay / lie 全家福：动词辨析的终点站}
\textbf{layer} 出自动词 \textbf{lay}（放置，及物：lay--laid--laid--laying）与 \textbf{lie}（躺；位于，不及物：lie--lay--lain--lying）构成英语最著名的混淆矩阵，另加“说谎”\textbf{lie}（lie--lied--lied--lying）三方混战：
\begin{center}
\begin{tabular}{lllll}
\toprule
原形 & 过去式 & 过去分词 & 现在分词 & 词义 \\
\midrule
lay & laid & laid & laying & 放置（及物） \\
lie & lay & lain & lying & 躺；位于（不及物） \\
lie & lied & lied & lying & 说谎（不及物） \\
\bottomrule
\end{tabular}
\end{center}
考点两连发：（1）\textit{She \textbf{lay} on the bed.}（躺的过去式，非“放置”）；（2）\textit{The book \textbf{was lying} open.}（现在分词）。判断口诀：有宾语用 lay 家族，无宾语用 lie 家族。
\end{gram}

\vw{legend}{/ˈledʒənd/}{n.}

\sense{n.}{传说，传奇故事}
\col{according to legend（据传说，叙事高频开头）}
\ex{According to legend, the lake was formed by a dragon's tear.}
\col{Legend has it that\dots（相传……）}
\ex{Legend has it that the castle's gold is still buried below.}

\sense{n.}{传奇人物}
\ex{The goalkeeper became a club legend after that final.}
\col{a living legend（在世传奇）}
\ex{At eighty, she is a living legend of Chinese cinema.}

\begin{lex}{legend / myth / folktale}
民俗学与文学理论的经典三分：\textbf{myth}（神话）——涉神灵、解释世界起源，被视为“真实”的神圣叙事；\textbf{legend}（传说）——涉人间英雄与史事，\textbf{似真}的世俗叙事（亚瑟王、白蛇传）；\textbf{folktale}（民间故事）——纯虚构娱乐，不pretend真实（灰姑娘）。口语里 legend 却滑向“传奇$=$了不起”，\textit{you legend!}（英俚“你真是个狠人！”）是词义褒扬的最新一站。
\end{lex}

\begin{origin}
经由古法语 \textit{legende} $\leftarrow$ 中世纪拉丁语 \textit{legenda}“\textbf{应被阅读之物}”——\textit{legere}“读”（原始印欧语 *\textit{leg-}“收集、拣选”$\to$把字句\textbf{拣}在一起即“读”）的动形词阴性复数。最初专指\textbf{圣徒传}：修道院餐堂里朗读的圣人生平（“读物”）。圣徒传渐染神异色彩，遂有“传说”义。同根词阵容豪华：\textbf{lecture}（读本$\to$讲座）、\textbf{legible}（可读的）、\textbf{intellectual}（inter-leg 拣选于内$\to$智识的）、\textbf{select}（se- 拣出$\to$挑选）、\textbf{elegant}（e- 拣出的$\to$精选的$\to$优雅的）、\textbf{diligent}?（否， diligent 出自 \textit{diligere}“珍爱挑选”，竟也算远亲）。一个“读”字，半部中世纪手抄本文化史。
\end{origin}

\vw{liberation}{/ˌlɪbəˈreɪʃn/}{n.\ \&\ liberty /ˈlɪbəti/ n.}

\sense{n.}{解放，摆脱}
\col{the liberation of / from sth}
\ex{The anniversary marks the liberation of the city in 1948.}
\col{women's liberation（妇女解放运动）}
\ex{Her grandmother was active in the women's liberation movement.}
\col{liberate / liberated（动词／形容词：被解放的；不受束缚的）}
\ex{Computers liberated office workers from mountains of paperwork.}

\sense{liberty n.}{自由}
\col{the Statue of Liberty（自由女神像）}
\ex{Millions of immigrants first saw the Statue of Liberty on arrival.}
\col{at liberty to do sth（有权做，正式）}
\ex{I'm not at liberty to discuss the details.}
\col{take the liberty of doing sth（冒昧做某事，礼貌套语）}
\ex{I took the liberty of reserving you a seat.}

\begin{gram}{lib / liber“自由”词族}
拉丁语 \textit{liber}“自由的”（原始印欧语 *\textit{leudh-}“成长、属于民众”）家族整齐：\textbf{liberty}（自由）、\textbf{liberate}（使自由）、\textbf{liberal}（自由的$\to$开明的；慷慨的；自由主义的；文科的 \textit{liberal arts}“自由人应受之艺”）、\textbf{libertine}（浪荡子“滥用自由者”）。历史彩蛋：古罗马对被打败的家长轻抚头部并\textit{mittam liberaliter}？更可靠的掌故是——被打 FREED 的奴隶戴\textit{pileus}自由帽；而英语 \textbf{free}（$<$原始日耳曼语，本义“亲爱的；不在奴籍的族人”）与拉丁 \textit{liber} 殊途同归于“自由”，朋友与自由同词（fri 和 free 同源：自由人即“自己人”）。
\end{gram}

\vw{like-minded}{/ˌlaɪk ˈmaɪndɪd/}{adj.}

\sense{adj.}{想法相同的，志趣相投的}
\col{like-minded people / friends / groups}
\ex{The club brings together like-minded lovers of classical music.}
\ex{Online forums help like-minded people find one another.}

\begin{gram}{复合形容词的“三层拼法”}
表语与定语的拼写差异是拼写题常客：（1）\textbf{like-minded}（连字符 + 名词过去分词，作\textbf{定语}）；（2）\textit{well-known} 定语 vs.\ \textit{The singer is well known.}（表语常断开）；（3）\textit{a five-year-old boy} vs.\ \textit{The boy is five years old}（定语单数无 -s，表语复数有 -s）。同构复合词：\textbf{narrow-minded}（心胸窄的）、\textbf{open-minded}（开明的）、\textbf{absent-minded}（心不在焉的）、\textbf{strong-minded}（意志坚定的）——\textit{mind} 是这批词的公共词根，mind$+$身体器官的搭配另见 \textit{make up one's mind / keep an open mind}。
\end{gram}

\vw{link}{/lɪŋk/}{vt.\ \&\ n.}

\sense{vt.}{联系，相关联；把……连接起来}
\col{link A with/to B}
\ex{The new bridge links the island with the mainland.}
\ex{Researchers have linked the disease to a common gene.}
\col{be linked to（与……有关联，新闻高频被动）}
\ex{Lack of sleep is directly linked to poor concentration.}

\sense{n.}{联系，连接；关系，纽带；链接}
\col{a link between A and B}
\ex{There is a close link between diet and health.}
\col{a weak / missing link（薄弱／缺失环节）}
\ex{The missing link in the theory was finally found in the fossils.}
\col{a video / web link（链接）}
\ex{Click the link below to watch the full interview.}

\begin{origin}
古英语 \textit{hlence}“\textbf{锁子甲的环}”（原始日耳曼语 *\textit{hlankō}，与 \textbf{lank}“细长的”？其近亲为 \textit{hlinkan}“弯曲”）：锁子甲由数千铁环\textbf{环环相扣}，“link”即一环。“关系、纽带”是 16 世纪的链条隐喻；生物学的“食物链”\textbf{food chain} 用了 chain，而“缺失环节”\textbf{missing link} 特指进化论中人与猿之间的化石缺环，19 世纪因此词而举世皆知。互联网的 hyperlink（超链接）让这个千年铁环词汇在数字时代重新上岗。
\end{origin}

\vw{logical}{/ˈlɒdʒɪkl/}{adj.}

\sense{adj.}{符合逻辑的；必然的，合乎情理的}
\col{a logical conclusion / step / explanation}
\ex{The logical conclusion is that the two events are unrelated.}
\ex{It's logical to assume prices will rise as demand grows.}
\col{logically（合逻辑地）}
\ex{Logically, the plan should work, but human nature intervened.}

\begin{origin}
经由拉丁语 \textit{logicus} $\leftarrow$ 希腊语 \textit{logikos}“\textbf{与推理相关的}”，出自 \textit{logos}“\textbf{言说、理性、道}”——希腊思想最重的词，原始印欧语 *\textit{leg-}“收集、言说”（与 legend 词条的拉丁 \textit{legere} 同源！希腊、拉丁两支都从“拣拾”走向“言说”）。\textit{logos} 家族在英语中构成半部学术史：\textbf{logic}（逻辑）、\textbf{biology / psychology / geology}（-logy“……之学”）、\textbf{dialogue}（dia- 对穿$\to$对话）、\textbf{apology}（apo- 挡开$\to$辩护（柏拉图《申辩篇》Apologia）$\to$道歉）、\textbf{analogy}（ana- 按比例$\to$类比）、\textbf{logo}（标识“商家的言语”）。《约翰福音》开篇 \textit{In the beginning was the Logos}（太初有道），和合本译作“道”——logos 与“道”，两大文明的最高理性词在此相遇。
\end{origin}

\vw{loose}{/luːs/}{adj.}

\sense{adj.}{不受约束的；未固定牢的；零散的；宽松的；疏松的}
\col{loose clothes / teeth / soil（宽松的衣物／松动的牙／松土）}
\ex{Wear loose, comfortable clothes for the long flight.}
\col{break / get loose（挣脱）}
\ex{One of the dogs broke loose and ran across the field.}
\col{let sth loose / let loose（释放；放任）}
\ex{The army let loose a barrage of shells.}
\ex{He let loose a roar of laughter.}
\col{loose change（零钱）；a loose end（未了结的琐事）}
\ex{He paid with loose change from his pocket.}
\ex{After the move, we still had several loose ends to tie up.}

\begin{gram}{loose / loosen / lose：一字母之差的三兄弟}
\textbf{loose} /luːs/ adj.\ 松的；\textbf{loosen} /ˈluːsn/ $v.$ 使松（\textit{loosen the knot}）；\textbf{lose} /luːz/ $v.$ 丢失（lose--lost--lost）。发音与拼写双重陷阱：\ding{55} \textit{I don't want to loose my keys.}（应为 lose）。动词化后缀 \textbf{-en}（loosen, widen, sharpen, strengthen）是“形容词转动词”的常规通道，一并记牢。反义链：loose $\leftrightarrow$ tight；lose $\leftrightarrow$ keep / find。
\end{gram}

\vw{lower}{/ˈləʊər/}{vt.\ \&\ adj.}

\sense{vt.}{降低，减少；把……放低，使……降下}
\col{lower prices / costs / one's voice}
\ex{The supermarket lowered prices to attract more shoppers.}
\ex{He lowered his voice to a whisper.}
\col{lower the flag（降旗）／ lower oneself（自贬）}
\ex{He refused to lower himself to their level of pettiness.}

\sense{adj.}{下面的；在底部的；低洼的}
\col{lower lip / floor / reaches（of a river）}
\ex{The floods hit the lower reaches of the river first.}
\ex{My office is on the lower floor, next to the library.}

\begin{gram}{“形容词 $+$ -er”的动词化通道}
\textbf{lower} $=$ low 的比较级 \textit{lower} 固化为\textbf{及物动词}“使降低”——同类：\textbf{further}（进一步$v.$）、\textbf{better}（改善$v.$：\textit{better one's record}）。但多数 -er 动词来自名词：\textit{hammer}（锤$v.$）、\textit{bottle}（装瓶）。判断捷径：lower 的动词义带宾语且常表“人为调控”（lower the volume），而“较低的”比较级义多用于比较结构（\textit{a lower price}）。反义动词 \textbf{raise / heighten / elevate} 与 lower 构成高低动词对。
\end{gram}

\vw{lyric}{/ˈlɪrɪk/}{n.\ \&\ adj.}

\sense{n.}{歌词；抒情诗}
\col{lyrics（复数：歌词）}
\ex{The lyrics of the song describe his hometown in spring.}
\col{a lyric poet（抒情诗人）}
\ex{Wordsworth is remembered as the greatest English lyric poet.}

\sense{adj.}{抒情的；吟唱的}
\col{lyric poetry / lyricism}
\ex{The film's ending has a lyric beauty that words can hardly describe.}

\begin{origin}
经由拉丁语 \textit{lyricus} $\leftarrow$ 希腊语 \textit{lyrikos}“\textbf{里拉琴的、随琴而歌}”，出自 \textit{lyra}“\textbf{里拉琴}”（lyre）——赫尔墨斯用龟壳发明、阿波罗执掌的\textbf{七弦琴}，古希腊抒情诗由它伴唱，故名“琴歌”。lyric poetry 在 19 世纪浪漫主义中达于极盛；爵士乐时代它被借指\textbf{流行歌词}（lyrics of a song），完成从缪斯到流行唱片的跨界。同根科学词 \textbf{lyrebird}（琴鸟，尾羽似琴）；而 \textbf{lyrical} 的口语化“热情洋溢的”（\textit{lyrical about her new job}）则是英国人少见的抒情时刻。
\end{origin}

\end{document}
